% !TeX root = ../main.tex

\begin{abstract}

WebAssembly (Wasm) 已廣泛部署於邊緣運算與無伺服器運算環境，這類場景對於冷啟動延遲特別敏感，因為其直接影響使用者感受到的回應速度。WasmEdge 為當前主流的伺服器端 Wasm 執行時期，目前僅以 LLVM 作為其唯一的 AOT 與 JIT 編譯器。儘管 LLVM 能產生高度優化的機器碼，其編譯流程具有顯著的延遲與記憶體開銷，並不適合資源受限的邊緣裝置或生命週期短的函式。WasmEdge 因此存在結構性的冷啟動瓶頸：每次啟動皆須先承擔完整的 LLVM 編譯成本，方能執行任何程式碼。

本論文提出一套適用於 WasmEdge 的分層編譯架構，目的在於將啟動延遲與最終程式碼品質解耦。本架構由三個元件組成。第一層為基線 JIT 編譯層 (Tier-1)，整合 dstogov/ir SSA 函式庫，於模組載入時以低成本將所有函式編譯為原生程式碼，立即提供執行能力。第二層為優化編譯層 (Tier-2)，重用 WasmEdge 既有的 LLVM 前端，透過合成 mini-module、三類 ABI thunk 橋接機制，以及納入靜態熱點呼叫對象的批次組合策略，選擇性地重新編譯熱點函式。第三項為在堆疊替換 (On-Stack Replacement, OSR) 機制，於最外層迴圈反向邊插入監測點，將正在執行中的 tier-1 框架於迴圈執行過程中遷移至 tier-2，解決函式入口替換在「單次呼叫者」場景下無法加速的問題。

本架構於 WasmEdge 中實作，並以 sightglass 測試套件進行評估。Tier-2 相較 tier-1 基線取得幾何平均 1.14 倍加速，達整體模組 LLVM-JIT 吞吐量的約百分之九十三；OSR 在一次性呼叫者測試中可將 tier-2 與 LLVM-JIT 之差距縮小至百分之二以內。在啟用所有層級的設定下，三十三個 sightglass-strong 基準測試於 O2 等級皆通過驗證。

\end{abstract}

\begin{abstract*}

WebAssembly (Wasm) is increasingly deployed in edge computing and serverless environments, where cold-start latency directly determines user-perceived responsiveness. WasmEdge, a leading server-side Wasm runtime, relies solely on LLVM for both ahead-of-time (AOT) and just-in-time (JIT) compilation. Although LLVM produces highly optimized machine code, its compilation pipeline incurs latency and memory overheads that are unsuitable for resource-constrained edge devices and short-lived functions. The runtime therefore exhibits a structural cold-start bottleneck: every invocation pays the full LLVM compile cost before any code executes.

This thesis presents a tiered compilation architecture for WasmEdge that decouples startup latency from peak code quality. The architecture comprises three components. First, a baseline JIT tier (Tier-1) integrates the dstogov/ir SSA library to compile every function eagerly at low cost, providing immediate execution capability at module load time. Second, an optimizing tier (Tier-2) reuses the existing WasmEdge LLVM frontend to selectively recompile hot functions through mini-module synthesis, three ABI bridge thunks, and a batch composition policy that includes statically-hot direct callees. Third, an on-stack replacement (OSR) mechanism instruments outermost-loop back-edges in tier-1 code and migrates currently-running tier-1 frames into tier-2 mid-execution, addressing the one-shot-caller scenario in which function-entry promotion alone delivers no speedup.

The architecture is implemented in WasmEdge and evaluated on the sightglass benchmark suite. Tier-2 promotion achieves a geometric-mean speedup of 1.14$\times$ over the tier-1 baseline and reaches approximately 93\% of whole-module LLVM-JIT throughput, while OSR closes the steady-state gap on one-shot-caller workloads to within 2\% of LLVM-JIT on representative kernels. All thirty-three sightglass-strong benchmark kernels pass at O2 under the combined configuration.

\end{abstract*}
